\chapter{Colectomy}

\section*{Introduction \& Clinical Indications}
Colectomy is a major surgical procedure involving the partial or complete removal of the large intestine, commonly known as the colon. This intervention is typically performed to treat a wide array of severe diseases affecting the colon, including malignant and benign neoplasms, chronic inflammatory conditions such as ulcerative colitis and Crohn's disease, complicated diverticular disease, and acute conditions like bowel obstruction or ischemia. The specific segment of the colon removed depends on the extent and location of the disease, ranging from a small segment (segmental colectomy) to the entire colon (total colectomy) and sometimes including the rectum (proctocolectomy). It is a definitive treatment aimed at curing disease, alleviating debilitating symptoms, or preventing life-threatening complications, significantly impacting patient quality of life and long-term health.

\begin{itemize}
  \item Bowel resection
  \item Colon resection
  \item Segmental colectomy
  \item Hemicolectomy (right or left)
  \item Total colectomy
  \item Proctocolectomy
\end{itemize}

The fundamental principle of colectomy is the surgical excision of diseased colonic tissue to eliminate the source of pathology, restore physiological function, or prevent further complications. For oncologic indications, the primary goal is complete tumor removal with adequate surgical margins and regional lymphadenectomy to achieve cure and prevent recurrence. This involves ligating the feeding vessels at their origin to ensure removal of the associated lymphatic basin. In inflammatory bowel disease, the aim is to resect the severely affected segment that is refractory to medical therapy, thereby resolving symptoms, preventing complications like perforation or toxic megacolon, and improving quality of life. For diverticular disease, resection targets the segment prone to recurrent inflammation or complications, typically the sigmoid colon.

Following resection, the remaining healthy ends of the colon are typically reconnected to restore bowel continuity, a procedure known as anastomosis. This can be performed using sutures or surgical stapling devices, aiming for a tension-free, well-perfused connection. In situations where an anastomosis is not safe or feasible due to significant inflammation, infection, gross contamination, or patient instability (e.g., sepsis), a stoma (colostomy or ileostomy) is created. This involves bringing the proximal bowel end through the abdominal wall to divert fecal contents into an external collection bag, either temporarily or permanently. The choice between anastomosis and stoma formation is a critical intraoperative decision based on the patient's clinical status, the quality of the bowel tissue, and the presence of contamination.

The concept of resecting portions of the colon dates back to the 19th century, with early attempts often fraught with high mortality rates due to infection and lack of effective anesthesia. The first successful elective colectomy is often attributed to Jean-François Reybard in 1833 for a rectal cancer, though his patient died shortly after. Theodor Billroth performed one of the earliest successful segmental colectomies in 1879. Significant advancements in surgical technique, particularly the adoption of aseptic principles by surgeons like Joseph Lister, and the development of safer anesthesia methods in the late 19th and early 20th centuries, dramatically improved outcomes.

The introduction of surgical stapling devices in the mid-20th century revolutionized anastomotic techniques, making them safer and more efficient by reducing operative time and improving consistency. The most significant modern evolution was the advent of laparoscopic colectomy in 1991 by Jacobs et al., which offered reduced pain, shorter hospital stays, and faster recovery compared to traditional open surgery. This minimally invasive approach has become the standard for many elective colectomies. Robotic-assisted colectomy, introduced in the early 2000s, further refines the minimally invasive approach, providing enhanced dexterity, tremor filtration, and 3D visualization, particularly beneficial for complex pelvic dissections.

\begin{itemize}
  \item To treat colon adenocarcinoma at various stages, aiming for curative resection with adequate oncologic margins and lymphadenectomy.
  \item To manage severe, complicated diverticulitis, especially with recurrent episodes, abscess formation, fistula, stricture, or perforation.
  \item To address medically refractory ulcerative colitis or Crohn's disease with complications such as strictures, fistulas, toxic megacolon, or high-grade dysplasia.
  \item To relieve acute or chronic large bowel obstruction caused by tumors, strictures, volvulus, or intussusception.
  \item To remove large, unresectable benign polyps or polyposis syndromes (e.g., Familial Adenomatous Polyposis (FAP), Lynch syndrome) with high malignant potential, often as a prophylactic measure.
  \item To treat severe colonic ischemia or gangrene requiring removal of necrotic bowel segments to prevent sepsis and multi-organ failure.
  \item To manage severe colonic trauma resulting in perforation, irreparable damage, or significant contamination.
  \item To treat colonic volvulus (e.g., sigmoid or cecal volvulus) that is recurrent or complicated by ischemia/perforation.
\end{itemize}

Colectomy serves as a primary, definitive treatment in many clinical scenarios. For localized colon cancer, it is the cornerstone of curative therapy, often followed by adjuvant chemotherapy depending on the stage and pathological findings. Similarly, for medically refractory ulcerative colitis, total proctocolectomy is considered a primary curative procedure, eliminating the risk of colorectal cancer. In cases of acute bowel obstruction, perforation, or severe ischemia, it is a primary, often emergent, life-saving intervention.

However, colectomy can also be a secondary or salvage procedure. For instance, in Crohn's disease, surgery is often reserved for complications that fail medical management, such as strictures, fistulas, or abscesses, rather than as a primary cure for the disease itself. It may also be a secondary procedure after failed endoscopic interventions for large polyps or as part of a multi-stage approach for complex inflammatory conditions or cancer with initial palliation (e.g., stoma creation for obstruction followed by definitive resection). The decision for colectomy is always made after careful consideration of the patient's overall health, the specific pathology, and the availability of less invasive alternatives.

\section*{Relevant Surgical Anatomy}
The large intestine, or colon, is approximately 1.5 meters long and extends from the ileocecal valve to the anus. It is divided into the cecum, ascending colon, transverse colon, descending colon, sigmoid colon, and rectum. The ascending and descending colon are retroperitoneal, while the transverse and sigmoid colon are intraperitoneal, suspended by mesenteries (transverse mesocolon and sigmoid mesocolon, respectively). These mesenteries contain the vascular, lymphatic, and nervous supply to the respective segments.

The arterial supply to the colon is derived from branches of the superior mesenteric artery (SMA) and inferior mesenteric artery (IMA). The SMA supplies the right colon (cecum, ascending colon, proximal two-thirds of the transverse colon) via the ileocolic, right colic, and middle colic arteries. The IMA supplies the left colon (distal one-third of the transverse colon, descending colon, sigmoid colon, and superior rectum) via the left colic, sigmoid, and superior rectal arteries. These arterial systems anastomose via the marginal artery of Drummond, which runs along the mesenteric border of the colon, providing collateral circulation. Venous drainage generally parallels the arterial supply, emptying into the superior mesenteric vein (SMV) and inferior mesenteric vein (IMV), which ultimately drain into the portal venous system. Lymphatic drainage follows the arterial supply, with lymph nodes located along the vessels (paracolic, intermediate, and principal nodes) draining into the superior and inferior mesenteric lymph node chains. The autonomic nervous system innervates the colon, with sympathetic fibers from the superior and inferior mesenteric plexuses and parasympathetic fibers from the vagus nerve (right colon) and pelvic splanchnic nerves (left colon and rectum). Key surgical landmarks include the ileocecal valve, the taeniae coli (longitudinal muscle bands that converge at the appendix and rectum), and the flexures (hepatic and splenic). Careful identification and preservation of the ureters, gonadal vessels, and major retroperitoneal structures are paramount during mobilization, especially in the left colon and pelvis.

\section*{Preoperative Workup \& Preparation}
Thorough preoperative preparation is crucial to optimize patient outcomes and minimize complications. This typically involves a comprehensive evaluation and specific interventions:
\begin{itemize}
  \item \textbf{Medical Evaluation:} A detailed history and physical examination, complete blood count, electrolytes, renal and liver function tests, coagulation profile, urinalysis, electrocardiogram, and often a chest X-ray are standard. For older patients or those with comorbidities, cardiac stress testing or pulmonary function tests may be required.
  \item \textbf{Anesthesia Clearance:} Evaluation by an anesthesiologist to assess fitness for general anesthesia and discuss pain management strategies (e.g., epidural, patient-controlled analgesia).
  \item \textbf{Imaging Studies:} Preoperative imaging such as CT scan of the abdomen and pelvis is essential to delineate the extent of disease, identify any metastases (for cancer), and plan the surgical approach. Colonoscopy is typically performed prior to elective colectomy to confirm diagnosis and rule out synchronous lesions.
  \item \textbf{Bowel Preparation:} For most elective colectomies, a mechanical bowel preparation (e.g., polyethylene glycol solution) is administered the day before surgery to cleanse the colon, often combined with oral antibiotics (e.g., neomycin and metronidazole) to reduce bacterial load and surgical site infection risk.
  \item \textbf{Medication Adjustment:} Patients on anticoagulants (e.g., warfarin, direct oral anticoagulants) or antiplatelet agents (e.g., aspirin, clopidogrel) will need these medications adjusted or temporarily stopped prior to surgery, often with bridging therapy. Diabetic patients require careful glucose management.
  \item \textbf{NPO Status:} Patients are instructed to be NPO (nil per os, nothing by mouth) for at least 6-8 hours before surgery to prevent aspiration during anesthesia.
  \item \textbf{Prophylactic Antibiotics:} Intravenous broad-spectrum antibiotics are administered within one hour prior to incision to reduce the risk of surgical site infection.
  \item \textbf{Informed Consent:} Detailed discussion with the surgeon regarding the procedure, potential risks, benefits, and alternatives, followed by obtaining written informed consent. This includes discussion of potential stoma formation.
\end{itemize}

\subsection*{Absolute Contraindications}
\begin{itemize}
  \item Uncorrectable severe coagulopathy or uncontrolled bleeding diathesis, posing an unacceptable risk of hemorrhage.
  \item End-stage multi-organ failure with a very poor prognosis and limited life expectancy, where the risks of surgery far outweigh any potential benefit.
  \item Diffuse metastatic disease with a very poor performance status (e.g., ECOG 3-4) and no palliative benefit from surgery.
  \item Patient refusal after thorough discussion of risks, benefits, and alternatives.
\end{itemize}

\subsection*{Relative Contraindications \& Precautions}
\begin{itemize}
  \item Severe cardiac or pulmonary disease that significantly increases anesthetic and surgical risk; requires extensive preoperative optimization and careful risk-benefit assessment.
  \item Active systemic infection or sepsis not originating from the colon, which should ideally be treated and resolved before elective surgery.
  \item Severe malnutrition or cachexia, which can impair wound healing, increase infection rates, and prolong recovery; may require preoperative nutritional support (e.g., TPN).
  \item Extensive intra-abdominal adhesions from previous surgeries, increasing the technical difficulty and risk of bowel or organ injury during dissection.
  \item Immunosuppression (e.g., from corticosteroids, biologics, chemotherapy) which can increase infection risk, impair wound healing, and potentially compromise anastomotic integrity.
  \item Pregnancy, especially in the first trimester, where surgery is generally avoided unless emergent and life-saving for the mother.
  \item Morbid obesity, which can complicate surgical access, increase operative time, and elevate the risk of wound complications and DVT/PE.
\end{itemize}

\section*{Surgical Technique \& Procedure}
Colectomy is performed under general anesthesia. The patient is typically positioned supine, though a modified lithotomy position may be used for low anterior resections involving the rectum to facilitate access to the pelvis and perineum. For laparoscopic or robotic approaches, the patient may be placed in a steep Trendelenburg position to aid visualization. The surgical approach can be open (via a single large abdominal incision, usually midline), laparoscopic (using several small incisions for trocars and specialized instruments), or robotic-assisted (a variation of laparoscopy with robotic arms).

The general steps for a colectomy include:
\begin{enumerate}
  \item \textbf{Anesthesia and Positioning:} General endotracheal anesthesia is induced. The patient is positioned appropriately, often with sequential compression devices applied to the lower extremities for deep vein thrombosis prophylaxis. The abdomen is prepped and draped in a sterile fashion.
  \item \textbf{Abdominal Access and Exploration:} For open surgery, a midline laparotomy is typically performed. For minimally invasive approaches, pneumoperitoneum is established, and trocars are placed. A thorough exploration of the abdominal cavity is performed to confirm the diagnosis, assess the extent of the disease, and rule out metastatic spread (for cancer).
  \item \textbf{Mobilization of the Colon:} The affected segment of the colon is carefully mobilized from its retroperitoneal attachments and surrounding structures. This involves incising peritoneal reflections and dissecting along avascular planes, taking care to preserve adjacent organs (e.g., ureters, duodenum, spleen) and nerves.
  \item \textbf{Vascular Ligation and Lymphadenectomy:} The blood vessels supplying the diseased segment are identified, ligated, and divided. For cancer, this often involves a high ligation of the main feeding artery (e.g., middle colic for right hemicolectomy, inferior mesenteric for left hemicolectomy) to ensure adequate lymph node harvest and oncologic clearance.
  \item \textbf{Bowel Transection:} The colon is transected proximally and distally to the diseased segment, ensuring clear margins, especially in oncologic cases. Surgical staplers are commonly used for this step, providing efficient and secure division.
  \item \textbf{Anastomosis or Stoma Formation:} The remaining healthy bowel ends are either reconnected (anastomosis) to restore continuity or brought out through the abdominal wall as a stoma (colostomy or ileostomy). The anastomosis can be hand-sewn or stapled, side-to-side, end-to-side, or end-to-end, depending on surgeon preference and anatomical considerations.
  \item \textbf{Irrigation and Closure:} The abdominal cavity is irrigated with saline, and meticulous hemostasis is ensured. Drains may be placed near the anastomosis or resection site to monitor for leaks or fluid collections, though their routine use is debated. The abdominal incision(s) are then closed in layers, with careful attention to fascial closure to prevent incisional hernia.
\end{enumerate}
The specific steps and extent of resection vary significantly based on the type of colectomy (e.g., right hemicolectomy, left hemicolectomy, sigmoid colectomy, total colectomy, total proctocolectomy) and the underlying pathology.

\section*{Complications \& Risks}
As with any major abdominal surgery, colectomy carries inherent risks, which can be broadly categorized as general surgical risks and procedure-specific complications.
\begin{itemize}
  \item \textbf{General Surgical Complications:}
    \begin{itemize}
      \item \textbf{Bleeding:} Intraoperative or postoperative hemorrhage, potentially requiring blood transfusion or reoperation.
      \item \textbf{Infection:} Surgical site infection (SSI) of the incision (superficial or deep), or deeper intra-abdominal infection (e.g., abscess).
      \item \textbf{Anesthesia Complications:} Adverse reactions to medications, respiratory or cardiac complications (e.g., pneumonia, myocardial infarction), stroke, or, rarely, death.
      \item \textbf{Deep Vein Thrombosis (DVT) and Pulmonary Embolism (PE):} Blood clots forming in the legs that can travel to the lungs, potentially fatal. Prophylaxis with anticoagulants and compression devices is routinely used.
    \end{itemize}
  \item \textbf{Procedure-Specific Risks:}
    \begin{itemize}
      \item \textbf{Anastomotic Leak:} Leakage of bowel contents from the surgical connection, leading to peritonitis, abscess formation, sepsis, and potentially requiring reoperation and/or creation of a temporary or permanent stoma. This is one of the most feared and serious complications.
      \item \textbf{Intra-abdominal Abscess:} Collection of pus within the abdominal cavity, often requiring percutaneous or surgical drainage.
      \item \textbf{Postoperative Ileus:} Prolonged absence of normal bowel function after surgery, leading to nausea, vomiting, and abdominal distension, delaying discharge.
      \item \textbf{Small Bowel Obstruction:} Blockage of the small intestine, often due to adhesions forming after surgery, potentially requiring reoperation.
      \item \textbf{Wound Dehiscence or Incisional Hernia:} Breakdown of the surgical incision or development of an incisional hernia at the site of the incision in the long term.
      \item \textbf{Stoma Complications:} If a stoma is created, potential issues include retraction, prolapse, stenosis, skin irritation, parastomal hernia, or high output.
      \item \textbf{Injury to Adjacent Organs:} Accidental damage to the ureters, bladder, small bowel, duodenum, spleen, or major retroperitoneal blood vessels during dissection, requiring immediate repair.
      \item \textbf{Nerve Injury:} Damage to autonomic nerves, particularly during pelvic dissection for rectal resections, leading to bladder, sexual, or bowel dysfunction.
      \item \textbf{Fecal Incontinence or Altered Bowel Habits:} Especially after low anterior resections or total proctocolectomy with ileal pouch-anal anastomosis, patients may experience increased stool frequency, urgency, or difficulty controlling bowel movements (low anterior resection syndrome).
    \end{itemize}
\end{itemize}

\section*{Postoperative Care \& Recovery}
Immediately after colectomy, the patient is transferred to the Post-Anesthesia Care Unit (PACU) for close monitoring of vital signs, pain level, and recovery from anesthesia. Pain management is a priority, often involving patient-controlled analgesia (PCA) or epidural catheters. Early ambulation is strongly encouraged, often within hours of surgery, as part of Enhanced Recovery After Surgery (ERAS) protocols, which aim to accelerate recovery and reduce complications.

Over the first 24-48 hours, patients are monitored for return of bowel function, indicated by the passage of flatus or a bowel movement. Diet is typically advanced gradually from clear liquids to a regular diet as tolerated, often starting on postoperative day 1. Intravenous fluids are administered until oral intake is sufficient. Drains, if placed, are monitored for output and typically removed when drainage is minimal. Hospital stay usually ranges from 3 to 7 days, depending on the complexity of the surgery, the patient's recovery trajectory, and the absence of complications. Upon discharge, patients receive instructions on wound care, pain medication, activity restrictions (e.g., no heavy lifting for 6-8 weeks), and warning signs to watch for, such as fever, increasing abdominal pain, or wound redness. Long-term recovery involves a gradual return to full activity over several weeks to months, with potential adjustments to diet and bowel habits, and for stoma patients, adaptation to stoma care.

During the colectomy, the surgeon meticulously documents intraoperative findings, including the precise location and extent of the disease, any involvement of adjacent structures, the presence of lymphadenopathy, and the integrity of the anastomosis or stoma. This information is critical for understanding the immediate surgical outcome and for guiding postoperative management.

The most critical "result" of a colectomy is the pathology report from the resected specimen. This report provides definitive information about the nature of the disease. For cancer, it details the tumor type (e.g., adenocarcinoma), histological grade, depth of invasion (T stage), involvement of regional lymph nodes (N stage), presence of vascular or perineural invasion, and most importantly, the status of the surgical margins (whether the tumor was completely removed with a rim of healthy tissue). This information is crucial for accurate cancer staging (TNM staging) and guides decisions regarding adjuvant therapies like chemotherapy. For inflammatory bowel disease, the pathology report confirms the diagnosis (e.g., ulcerative colitis, Crohn's disease), assesses the severity and extent of inflammation, and identifies any dysplasia or malignant transformation. For diverticular disease, it confirms the presence of diverticula and inflammation.

A normal, successful postoperative course following colectomy signifies effective treatment of the underlying pathology and an uncomplicated recovery. Patients can expect a gradual resolution of pain, which is typically managed with oral analgesics after discharge. Bowel function should normalize over several days to weeks, with the passage of flatus and stool, allowing for a progressive return to a regular diet. The surgical incision(s) should heal cleanly without signs of infection, dehiscence, or excessive pain. For cancer patients, a successful course includes clear surgical margins and, ideally, negative lymph nodes, indicating a high likelihood of cure or long-term disease control. Patients should experience a progressive increase in energy and activity levels, with a return to most normal daily activities within 6-8 weeks, though full recovery can take several months. The ultimate goal is alleviation of preoperative symptoms and an improved quality of life.

Post-colectomy follow-up is essential for monitoring recovery, managing long-term effects, and detecting recurrence or new issues.
\begin{itemize}
  \item \textbf{Postoperative Clinic Visits:} Typically scheduled 1-2 weeks after discharge for wound check, staple/suture removal, and general assessment of recovery. Further visits may be scheduled based on individual needs and pathology results.
  \item \textbf{Stoma Care and Education:} If a stoma was created, patients will receive extensive education and support from an ostomy nurse for appliance management, skin care, dietary adjustments, and psychosocial support.
  \item \textbf{Oncology Follow-up:} For cancer patients, referral to a medical oncologist is standard to discuss adjuvant chemotherapy and long-term surveillance plans, which may include regular blood tests (e.g., CEA levels), CT scans of the chest/abdomen/pelvis, and surveillance colonoscopies.
  \item \textbf{Surveillance Colonoscopy:} Recommended at specific intervals (e.g., 1 year after surgery, then every 3-5 years) to monitor for recurrence, new polyps, or metachronous cancers in the remaining colon, especially for patients with a history of polyps or IBD.
  \item \textbf{Dietary and Lifestyle Counseling:} Guidance on diet, hydration, and activity levels to optimize bowel function and overall health, particularly for patients with altered bowel habits or a stoma.
  \item \textbf{Physical Therapy/Rehabilitation:} May be recommended for patients with significant deconditioning or specific functional deficits to aid in regaining strength and mobility.
\end{itemize}
Patients are advised to contact their surgeon immediately for warning signs such as persistent fever (over 101.5$^{\circ}$F or 38.6$^{\circ}$C), severe or worsening abdominal pain, persistent nausea/vomiting, inability to pass gas or stool, significant wound changes (redness, swelling, pus), or heavy rectal bleeding.

\section*{Outcomes \& Prognosis}
Colectomy is a frequently performed major abdominal surgery, with its incidence largely driven by the prevalence of colorectal cancer and inflammatory bowel disease. Colorectal cancer is the third most common cancer and the second leading cause of cancer-related death globally, with over 1.9 million new cases annually. The incidence of colectomy for cancer increases with age, peaking in individuals over 60, and shows slight variations by sex and ethnicity. Inflammatory bowel disease (Crohn's disease and ulcerative colitis) affects millions worldwide, with a significant proportion of patients (up to 25\% for UC, 75\% for CD) eventually requiring surgery due to medical treatment failure or complications. Complicated diverticular disease also contributes substantially to colectomy rates, particularly in older populations, with an estimated 15-25\% of patients with diverticulitis requiring surgery. The overall mortality rate for elective colectomy is generally low, ranging from 1-3\% in specialized centers, but can be significantly higher (up to 10-20\%) in emergent settings or for high-risk patients. Morbidity rates, encompassing complications like anastomotic leak, infection, and ileus, typically range from 15-30\%.

The outcomes and prognosis following colectomy are highly dependent on the underlying indication, the stage of the disease, and the patient's overall health. For localized colon cancer, colectomy offers a high chance of cure, with 5-year survival rates ranging from over 90\% for Stage I disease to 70-80\% for Stage II, and 50-70\% for Stage III, particularly when combined with adjuvant chemotherapy. Prognostic factors include tumor stage, lymph node involvement, surgical margin status, tumor biology, and patient comorbidities. Landmark clinical trials, such as the COLOR II trial, have demonstrated the non-inferiority of laparoscopic colectomy compared to open surgery for colon cancer in terms of oncologic outcomes. For medically refractory ulcerative colitis, total proctocolectomy is curative, eliminating the risk of colorectal cancer and significantly improving quality of life, although patients may experience altered bowel habits or pouch-related complications if an ileal pouch-anal anastomosis is performed. In Crohn's disease, colectomy can effectively manage complications and provide symptom relief, but it is not curative, and recurrence in the remaining bowel segments is common, necessitating ongoing medical management. For diverticular disease, elective colectomy significantly reduces the risk of recurrent diverticulitis. Overall, most patients experience a good functional recovery and improved quality of life after colectomy, though some may face long-term challenges such as altered bowel function, stoma management, or the need for ongoing surveillance.

\section*{References}
\begin{enumerate}
  \item Brunicardi FC, Andersen DK, Billiar TR, et al. Schwartz's Principles of Surgery. 11th ed. McGraw-Hill Education; 2019.
  \item American Society of Colon and Rectal Surgeons. Clinical Practice Guidelines for Colon Cancer. 2023. Available from: \url{https://fascrs.org/patients/diseases-and-conditions/colon-cancer}
  \item MedlinePlus. Colectomy. U.S. National Library of Medicine; 2024. Available from: \url{https://medlineplus.gov/colectomy.html}
  \item Buunen M, Veldkamp M, Hop WC, et al. Survival after laparoscopic surgery versus open surgery for colon cancer: long-term outcome of a randomised clinical trial. Lancet Oncol. 2009;10(1):44-51. (COLOR trial)
\end{enumerate}

\clearpage
